\documentclass[12pt,a4paper]{article}

\usepackage[utf8]{inputenc}
\usepackage[T1]{fontenc}
\usepackage[margin=2.2cm]{geometry}
\usepackage{setspace}
\usepackage{hyperref}
\usepackage{booktabs}
\usepackage{enumitem}
\setlist[itemize]{noitemsep,topsep=3pt}
\setlist[enumerate]{noitemsep,topsep=3pt}
\onehalfspacing

\title{PneumoScan -- Guide de Revision Personnel\\\large Version simple pour preparer les questions du jury (10 min)}
\author{Anis Aissani}
\date{\today}

\begin{document}
\maketitle
\tableofcontents
\newpage

\section{Objectif de ce document}
Ce document est fait pour \textbf{comprendre rapidement tout le projet} et pouvoir repondre clairement aux questions du jury.

Idee cle:
\begin{itemize}
  \item Ton projet n'est pas seulement un modele ML.
  \item C'est une \textbf{application complete}: backend, frontend, base de donnees, securite, deploiement Docker, tests.
\end{itemize}

\section{Le projet en 30 secondes}
\subsection{Probleme}
En radiologie, les radios peuvent etre lues en FIFO (ordre d'arrivee). Un cas grave peut donc attendre.

\subsection{Solution}
PneumoScan analyse une radio thoracique et retourne:
\begin{itemize}
  \item une prediction (\texttt{PNEUMONIA} ou \texttt{NORMAL}),
  \item un score de confiance,
  \item un niveau de triage (\texttt{CRITICAL}, \texttt{MODERATE}, \texttt{ROUTINE}),
  \item une heatmap explicative,
  \item une trace en base de donnees.
\end{itemize}

\subsection{Valeur principale a dire au jury}
\textbf{Tu proposes un prototype deployable et tracable}, pas juste un script d'IA.

\subsection{Repartition binome a dire clairement}
\begin{itemize}
  \item Mon binome a travaille sur l'entrainement et la validation du modele.
  \item J'ai travaille sur l'application complete (API FastAPI, interface, securite, BDD, Docker).
  \item Mon integration applicative s'appuie sur son travail modele, tout en gardant une vision globale de toute la chaine.
\end{itemize}

\subsection{UP adapte (a reciter tel quel)}
\begin{itemize}
  \item Inception (2 sem.) $\rightarrow$ Etude existant, choix stack, repartition binome.
  \item Elaboration (4 sem.) $\rightarrow$ UML, schema BDD, maquettes, binome entraine modele.
  \item Construction (8 sem.) $\rightarrow$ API FastAPI, Vue.js, integration modele, tests.
  \item Transition (2 sem.) $\rightarrow$ Docker, rapport, corrections finales.
\end{itemize}

\section{Architecture globale (simple)}
\subsection{Blocs techniques}
\begin{itemize}
  \item \textbf{Frontend}: interface utilisateur (Vue.js dans \texttt{index.html}).
  \item \textbf{Backend}: API FastAPI (routes REST, auth, orchestration).
  \item \textbf{ML}: modele principal EfficientNet-B0 + SVM (fallback HOG + PCA + SVM).
  \item \textbf{BDD}: SQLite pour historiser predictions et sessions.
  \item \textbf{Docker Compose}: deploiement reproductible en 2 services (backend + frontend).
\end{itemize}

\subsection{Flux principal}
\begin{enumerate}
  \item Le radiologue se connecte (JWT).
  \item Il ouvre une session de travail.
  \item Il envoie une image a \texttt{/predict}.
  \item Backend valide fichier + garde OOD.
  \item Pipeline ML calcule prediction + probabilite.
  \item Backend calcule le triage, genere la heatmap et stocke en base.
  \item Frontend affiche resultat, triage et historique.
\end{enumerate}

\section{Backend: ce que chaque fichier fait}
\begin{table}[h!]
\centering
\begin{tabular}{p{4cm} p{10.5cm}}
\toprule
\textbf{Fichier} & \textbf{Role}\
\midrule
\texttt{main.py} & Routes HTTP FastAPI (controller layer), validation HTTP, codes d'erreur, appels aux services. \\
\texttt{services.py} & Orchestration metier (prediction, sessions, historique, dashboard). \\
\texttt{settings.py} & Configuration centralisee (version app, dossier heatmaps, etc.). \\
\texttt{auth.py} & JWT, RBAC, verification utilisateur, dependances \texttt{require\_radiologist}/\texttt{require\_admin}. \\
\texttt{database.py} & CRUD SQLite, triage, sessions, stats dashboard. \\
\texttt{predictor.py} & OOD guard, preprocess, extraction features (EfficientNet ou HOG fallback), infererence SVM, heatmap, anonymisation. \\
\texttt{utils.py} & Validation des fichiers uploades (type, extension, taille). \\
\texttt{scanner.py} & Traitement batch automatique des images dans \texttt{incoming\_scans}. \\
\texttt{test\_backend.py} & Tests unitaires (32 tests). \\
\bottomrule
\end{tabular}
\end{table}

\subsection{Phrase simple pour expliquer la refactorisation}
\textbf{Avant}, beaucoup de logique etait concentree dans les routes.\\
\textbf{Apres}, on separe:
\begin{itemize}
  \item HTTP (controllers),
  \item logique metier (services),
  \item acces donnees (database),
  \item IA (predictor),
  \item configuration (settings).
\end{itemize}
C'est plus maintenable, testable, et plus clair a presenter.

\section{Le pipeline IA, tres simplement}
\subsection{1) Validation OOD (anti faux positifs absurdes)}
Le modele ne doit pas traiter n'importe quelle image (photo couleur, document, etc.).

Verification OOD:
\begin{itemize}
  \item saturation HSV faible (une radio est quasi en niveaux de gris),
  \item densite de bords Canny raisonnable.
\end{itemize}
Si invalide: HTTP 400.

\subsection{2) Pretraitement}
\begin{itemize}
  \item conversion en niveaux de gris,
  \item mise au format 224x224,
  \item centrage sur canvas noir pour garder le ratio.
\end{itemize}

\subsection{3) Features}
Voie principale: embedding EfficientNet-B0 (1280 dimensions), puis eventuelle reduction selon l'artefact.\
Voie fallback: vecteur HOG (576 features).

\subsection{4) Classification SVM}
Le modele predit classe et probabilite de pneumonie.

\subsection{5) Triage}
Regles metier:
\begin{itemize}
  \item \texttt{CRITICAL}: pneumonie et $p > 0.85$,
  \item \texttt{MODERATE}: pneumonie et $0.15 \le p \le 0.85$,
  \item \texttt{ROUTINE}: normal, ou pneumonie a faible confiance $p < 0.15$.
\end{itemize}

\section{Securite (a retenir pour le jury)}
\begin{itemize}
  \item Authentification par JWT (duree 8h).
  \item RBAC avec 2 roles: \texttt{radiologist} et \texttt{admin}.
  \item Mots de passe haches (bcrypt/passlib).
  \item Sessions anti-zombie: fermeture auto des sessions ouvertes precedemment.
\end{itemize}

\section{Base de donnees (SQLite)}
Deux tables principales:
\begin{itemize}
  \item \texttt{predictions}: resultat d'analyse (label, score, triage, heatmap, operateur, session).
  \item \texttt{sessions}: session de travail (ouverte/fermee, nombre de patients, operateur).
\end{itemize}

Pourquoi SQLite ici:
\begin{itemize}
  \item simple et suffisant pour un prototype mono-instance,
  \item migration PostgreSQL possible ensuite.
\end{itemize}

\section{Frontend (ce que tu dois savoir expliquer)}
\begin{itemize}
  \item Login + stockage token en \texttt{sessionStorage}.
  \item Ouverture/cloture de session.
  \item Upload image + option anonymisation.
  \item Affichage prediction, score, triage, heatmap.
  \item Historique filtrable + export CSV.
  \item Dashboard admin (KPI + graphiques).
  \item Generation rapport PDF.
\end{itemize}

\section{Docker et deploiement (point tres important)}
\subsection{Pourquoi c'est central}
Docker prouve:
\begin{itemize}
  \item reproductibilite,
  \item portabilite,
  \item industrialisation minimale du prototype.
\end{itemize}

\subsection{Commandes essentielles}
\begin{itemize}
  \item Lancer: \texttt{docker compose up --build}
  \item En arriere-plan: \texttt{docker compose up -d --build}
  \item Etat: \texttt{docker compose ps}
  \item Logs: \texttt{docker compose logs -f}
  \item Arreter: \texttt{docker compose down}
\end{itemize}

\textbf{Point de demo:} l'interface est accessible sur \texttt{http://localhost:3000}.

\subsection{Attention pratique}
\texttt{docker-compose-v2} n'est pas la bonne commande ici.\\
Utiliser \texttt{docker compose ...} (plugin V2).

\section{Tests et qualite}
\begin{itemize}
  \item 32 tests unitaires, passes.
  \item Validation upload, preprocess, HOG, triage, anonymisation.
  \item Objectif: limiter les regressions apres refactorisation.
\end{itemize}

\section{Limites actuelles (important pour paraitre mature)}
\begin{itemize}
  \item Support DICOM présent, mais sans intégration PACS hospitalier dans cette version.
  \item Le mode fallback HOG+SVM est moins performant que le modele principal EfficientNet+SVM.
  \item Prototype non certifie (pas dispositif medical CE/FDA).
  \item SQLite non ideal pour charge multi-utilisateurs a grande echelle.
\end{itemize}

\section{Evolutions proposees}
\begin{itemize}
  \item Support DICOM + integration PACS.
  \item Fine-tuning d'EfficientNet-B0 sur des donnees plus diversifiees.
  \item Explicabilite avancee (Grad-CAM, SHAP).
  \item Migration PostgreSQL + CI/CD + monitoring.
\end{itemize}

\section{Questions probables du jury + reponses courtes}
\subsection{Q1. Pourquoi EfficientNet-B0 + SVM ?}
\textbf{Reponse courte:} c'est le meilleur compromis robuste dans les resultats du binome (Accuracy 86.2\%, F1 88.4\%). On evite de retenir la concat perfect-score (100\%) car elle peut indiquer une fuite de donnees.

\subsection{Q2. Comment eviter qu'une photo classique soit analysee ?}
\textbf{Reponse courte:} garde OOD avant infererence (saturation HSV + densite Canny). Si invalide, rejet HTTP 400.

\subsection{Q3. Pourquoi la refactorisation etait necessaire ?}
\textbf{Reponse courte:} pour separer les responsabilites (routes/services/data/ML/config), faciliter les tests et rendre l'architecture defendable au jury.

\subsection{Q4. Que garantit la securite ?}
\textbf{Reponse courte:} JWT + RBAC + mots de passe haches + sessions anti-zombie. Controle role-based sur les endpoints sensibles.

\subsection{Q5. Pourquoi Docker est un point majeur ?}
\textbf{Reponse courte:} car il transforme le projet en application deployable/reproductible, pas juste un code local. C'est le socle de l'industrialisation.

\subsection{Q6. Pourquoi SQLite et pas PostgreSQL ?}
\textbf{Reponse courte:} SQLite est adaptee au prototype mono-instance. PostgreSQL est prevu pour l'evolution multi-utilisateurs.

\subsection{Q7. Comment prouves-tu que ton code est fiable ?}
\textbf{Reponse courte:} suite de 32 tests unitaires passes, plus verification runtime via Docker Compose et endpoints fonctionnels.

\subsection{Q8. Limite la plus importante ?}
\textbf{Reponse courte:} absence de certification medicale et besoin de validation externe multicentrique avant usage clinique.

\section{Mini script oral (2 minutes) pour les questions)}
\begin{enumerate}
  \item Mon projet repond a un probleme de priorisation clinique en radiologie.
  \item J'ai construit une application complete: API, interface, BDD, securite, deployment Docker.
  \item Le coeur IA retenu est EfficientNet-B0+SVM, avec fallback HOG+SVM et garde OOD pour eviter les entrees invalides.
  \item La logique de triage transforme la prediction en priorite operationnelle.
  \item La refactorisation en couches rend le code propre, testable et maintenable.
  \item Docker Compose est central car il garantit un deploiement reproductible.
  \item J'assume les limites (prototype non certifie, pas d'integration PACS), et j'ai une roadmap claire.
\end{enumerate}

\section{Checklist de revision finale}
Avant la soutenance, verifier que tu sais expliquer sans hesiter:
\begin{itemize}
  \item le flux complet de \texttt{/predict},
  \item les seuils de triage,
  \item pourquoi OOD est indispensable,
  \item ce que Docker apporte concretement,
  \item les limites + evolutions.
\end{itemize}

\vspace{10pt}
\textbf{Conseil final:} pendant les questions, reponds en 3 temps: \textit{contexte -> choix -> impact}.\\
Exemple: \textit{"J'ai choisi X parce que Y, et l'impact est Z"}.

\end{document}
